% qss-odbc-manual.tex --- A short manual for qss-odbc
%
% Copyright (C) 2026 Friedrich Dominicus
% SPDX-License-Identifier: MIT
%
% Build with LuaLaTeX (twice):  lualatex qss-odbc-manual.tex

\documentclass[
  paper=a4, fontsize=11pt, parskip=half, headings=small,
  toc=bibliography, numbers=noenddot, DIV=12
]{scrreprt}

\usepackage{fontspec}
\setmonofont{DejaVu Sans Mono}[Scale=MatchLowercase]
\usepackage[english]{babel}
\usepackage{microtype}
\usepackage[autostyle]{csquotes}
\usepackage{xcolor}
\usepackage{booktabs, longtable, array, tabularx}
\usepackage{fancyvrb}
\usepackage{listings}
\usepackage[automark]{scrlayer-scrpage}
\usepackage[hidelinks]{hyperref}

\definecolor{qssframe}{gray}{0.75}
\definecolor{qssback}{gray}{0.96}

\newcommand{\code}[1]{\texttt{#1}}
\newcommand{\file}[1]{\texttt{#1}}
\newcommand{\kbd}[1]{\texttt{#1}}
\newcommand{\qss}{\textsf{qss}}

\fvset{fontsize=\small, frame=leftline, framerule=1.5pt, rulecolor=\color{qssframe},
       xleftmargin=1em, framesep=0.8em}

\lstdefinestyle{elisp}{
  language=Lisp, basicstyle=\ttfamily\small, columns=fullflexible, keepspaces=true,
  upquote=true, showstringspaces=false, frame=leftline, framerule=1.5pt,
  rulecolor=\color{qssframe}, xleftmargin=1em, framesep=0.8em,
  commentstyle=\color{black!60}, keywordstyle={}, ndkeywordstyle={},
}
\lstset{style=elisp}

\pagestyle{scrheadings}
\KOMAoptions{headsepline=true}

\title{qss-odbc}
\subtitle{ODBC databases from Emacs Lisp\\[1ex]
  \normalsize A short manual}
\author{Friedrich Dominicus}
\newcommand{\qssversion}{0.3.0}
\date{Version \qssversion}

\begin{document}

\maketitle
\tableofcontents

% ======================================================================
\chapter{What it is}

Any database with an ODBC driver --- PostgreSQL, MariaDB, SQL Server, Access,
SQLite \ldots{} --- from Emacs Lisp:

\begin{lstlisting}
(require 'qss-odbc)
(setq con (qss-odbc-connect "test"))      ; a DSN, or a connection string
(qss-odbc-query con "SELECT id, name FROM person WHERE id > ?" 10)
;; => (("id" "name") (11 "Müller") (12 "Ærø"))
(qss-odbc-do-query ((id name) con "SELECT id, name FROM person")
  (message "%s %s" id name))
(qss-odbc-with-transaction con
  (qss-odbc-execute con "UPDATE person SET name = ? WHERE id = ?" "X" 11))
(qss-odbc-disconnect con)
\end{lstlisting}

Emacs Lisp has no foreign function interface, and Emacs will not get one: it
would let any library into Emacs, free or not. What there is are dynamic
modules, written in C against \file{emacs-module.h}. So the calls into ODBC
live in a small module, \file{module/qss-odbc-module.c}, and
\file{qss-odbc.el} is the library on top of it.

The division of labour: \textbf{C does what needs memory and pointers} ---
handles, buffers, UTF-16, binding parameters, reading columns piece by piece,
collecting diagnostics --- and returns raw values. \textbf{Lisp interprets
them}: dates, exact decimals, booleans, the column descriptions. No pointer
reaches Lisp, so no mistake in Lisp can crash Emacs.

The rules --- which type is bound how, lengths in which unit, buffer sizes,
the quirks of drivers --- follow plain-odbc in its QSS version for Common
Lisp, tested against Access, SQLite, PostgreSQL and MariaDB.

\qss{} (\url{https://github.com/FDominicus/qss}) uses this package as a
source of tables: \code{(qss-> (odbc "test" "person") (where (> :id 10)))}.

% ======================================================================
\chapter{Building the module}

\begin{Verbatim}
M-x qss-odbc-build-module
\end{Verbatim}

It runs \kbd{make} in \file{module/}, shows the output in
\code{*qss-odbc build*} and loads the result. It needs a C compiler, the ODBC
headers and a driver manager:

\noindent\begin{tabularx}{\linewidth}{@{}lX@{}}
\toprule
Debian, Ubuntu & \kbd{apt install unixodbc-dev} (and a driver: \code{odbc-postgresql}, \code{odbc-mariadb}, \code{libsqliteodbc}) \\
macOS & unixODBC or iODBC; \code{ODBC=iodbc} for the latter \\
Windows & MinGW-w64 (MSYS2); the driver manager \file{odbc32} is built in \\
\bottomrule
\end{tabularx}

\medskip
If \file{emacs-module.h} is not installed (an Emacs built in its source
tree), set \code{qss-odbc-make-arguments} to
\code{("EMACS\_INCLUDE=/path/to/emacs/src")}. Emacs must be 29.1 or later and
built with modules (\code{module-file-suffix} is not nil). If the module is
missing, \code{qss-odbc-connect} offers to build it.

% ======================================================================
\chapter{Connecting and asking}

\section{Connections}

\code{(qss-odbc-connect DSN \&rest OPTIONS)}: DSN is a data source name
(\kbd{odbcinst -q -s}) or a full connection string. OPTIONS: \code{:user},
\code{:password}, \code{:autocommit} (default \code{t}), \code{:utf-8} (see
chapter~\ref{chap:drivers}). The connection prints with the password hidden:
\code{\#<qss-odbc PostgreSQL DSN=test;UID=frido;PWD=***>}.

Connections start with autocommit on: every statement is its own
transaction. \code{qss-odbc-with-transaction} runs a body in one, committed
at the end and rolled back on any non-local exit; \code{qss-odbc-commit} and
\code{qss-odbc-rollback} are there too.

\section{Queries and statements}

\noindent\begin{tabularx}{\linewidth}{@{}lX@{}}
\toprule
\code{qss-odbc-query} & rows as lists, the first the column names (as \code{sqlite-select} with \code{full}) \\
\code{qss-odbc-select} & \code{(COLUMNS . ROWS)}; columns as property lists \code{:name :type :size :scale :nullable :type-name} \\
\code{qss-odbc-select-all} & all result sets of a statement \\
\code{qss-odbc-execute} & a statement; the number of rows changed \\
\code{qss-odbc-map-query} & a function on each row, one row in memory at a time \\
\code{qss-odbc-do-query} & the same as a loop; \code{cl-return} leaves it \\
\code{qss-odbc-tables}, \code{-table-columns}, \code{-primary-key} & the catalog: SQLTables, SQLColumns, SQLPrimaryKeys; the last two take a schema \\
\bottomrule
\end{tabularx}

\medskip
Values go in as parameters (the \code{?} in the SQL), never as text in the
SQL. A statement the garbage collector finds unreferenced is freed; closing a
connection closes its statements.

\section{Errors}

Errors are signalled as \code{qss-odbc-error} with the data
\code{(MESSAGE SQLSTATE NATIVE-ERROR DIAGNOSTICS)}; DIAGNOSTICS holds every
record. \code{qss-odbc-error-sqlstate} and \code{-message} take them apart:

\begin{lstlisting}
(condition-case err (qss-odbc-query con "SELECT * FROM nothing")
  (qss-odbc-error (qss-odbc-error-sqlstate err)))   ; => "42P01"
\end{lstlisting}

% ======================================================================
\chapter{Values}

\noindent\begin{tabularx}{\linewidth}{@{}lX@{}}
\toprule
\textbf{From the database} & \textbf{In Lisp} \\
\midrule
NULL & \code{nil} \\
integers & integers (bignums if need be) \\
REAL, FLOAT, DOUBLE & floats \\
NUMERIC(p,s), DECIMAL & integers scaled by $10^s$: 1094.80 in a NUMERIC(10,2) is \code{109480}; \code{qss-odbc-decimal-string} shows it \\
BIT, BOOLEAN & \code{t}, \code{nil} \\
DATE, TIME, TIMESTAMP & text, the wall clock as stored: \code{"2026-10-30 09:56:00.123456"}; \code{qss-odbc-parse-time} makes a Lisp time \\
binary & unibyte strings \\
text, everything else & strings \\
\bottomrule
\end{tabularx}

\medskip
Decimals are exact: a currency amount is never a float on the way. The
scale is in the column (\code{qss-odbc-column-scale}), so 109480 and 2 are the
amount; a column of a fixed scale --- Access currency has 4 --- adds and
compares as integers.

\begin{tabularx}{\linewidth}{@{}lX@{}}
\toprule
\textbf{Parameter} & \textbf{Bound as} \\
\midrule
\code{nil}, \code{t} & NULL, TRUE; FALSE is \code{(qss-odbc-bool nil)} \\
integer, float, string & INTEGER (beyond 32 bit NUMERIC), DOUBLE, WVARCHAR \\
\code{(qss-odbc-decimal 109480 2)} & NUMERIC 1094.80 \\
\code{(qss-odbc-binary BYTES)} & VARBINARY \\
\code{(qss-odbc-date "2026-10-30")} & DATE \\
\code{(qss-odbc-timestamp "2026-10-30 09:56")} & TIMESTAMP (or a Lisp time, taken in UTC) \\
\bottomrule
\end{tabularx}

% ======================================================================
\chapter{Drivers}\label{chap:drivers}

Unicode drivers get UTF-16, through the W functions of ODBC. ANSI drivers that
expect UTF-8 --- SQLite, and psqlODBC's ANSI driver, often registered by
mistake under the name \emph{PostgreSQL Unicode} --- get SQL text,
parameters, character columns and messages as UTF-8 bytes. Through the W
functions unixODBC would turn them into Latin-1 (SQLite lost even
\emph{ü}), and psqlODBC ANSI rejects UTF-16 parameters with HYC00
\emph{Unrecognized C\_parameter type}. Which drivers count as such:
\code{qss-odbc-ansi-utf-8-drivers}, or \code{:utf-8 t} at
\code{qss-odbc-connect}.

What the tests found:

\begin{description}
\item[SQLite] reports VARCHAR through SQLColAttribute for BIGINT and BLOB
  columns (the type comes from SQLDescribeCol); keeps 15 digits of a double;
  has no exact numbers.
\item[psqlODBC] reports booleans as VARCHAR \code{"1"}/\code{"0"} (the type
  name \code{bool} decides); splits long messages into several records at the
  size of the first request.
\item[MariaDB, MySQL] have no boolean: BOOLEAN is TINYINT(1) and comes as 1 or
  0. TEXT and BLOB end at 64\,KB, TIMESTAMP covers 1970 to 2038 (LONGTEXT,
  LONGBLOB, DATETIME hold more). Connector/ODBC before 3.1.24 and 3.2.10
  miscounts characters U+xxD8 to U+xxDB (\emph{Ø}). An UPDATE counts only rows
  whose values changed.
\end{description}

% ======================================================================
\chapter{Tests}

\begin{Verbatim}
QSS_ODBC_DSNS="test mytest" emacs -Q --batch -L . -l qss-odbc-test \
  -f ert-run-tests-batch-and-exit
\end{Verbatim}

Unit tests run always; database tests against each DSN in
\code{QSS\_ODBC\_DSNS} (default \code{test}), skipped with the reason if the
module is not built or a DSN cannot be reached. They create and drop the
table \code{qss\_odbc\_t}. Text in many scripts and with surrogate pairs on
the piece boundaries of SQLGetData, integers at 32 and 64 bit, exact
decimals, dates before 1900 and at both changes of summer time, binary data,
errors, transactions, 2000 rows one at a time, 3000 statements left to the
garbage collector, catalog functions.

Tested with PostgreSQL 16 (psqlODBC Unicode and ANSI, also a database that
sorts by German rules), MariaDB 10.11 (Connector/ODBC 3.1.15) and SQLite3
ODBC; the module under valgrind without errors or leaks.

% ======================================================================
\chapter{Files}

\noindent\begin{tabularx}{\linewidth}{@{}lX@{}}
\toprule
\file{qss-odbc.el} & the library: connections, queries, values, the catalog, building the module \\
\file{module/qss-odbc-module.c} & the C module: handles, UTF-16, parameters, columns, diagnostics \\
\file{module/Makefile} & Linux, macOS (iODBC), Windows (MinGW) \\
\file{qss-odbc-test.el} & the tests \\
\file{doc/} & this manual \\
\bottomrule
\end{tabularx}

\end{document}
